\section{Negative Results and Alternative Approaches}
\label{sec:5_8}

In this section, we document the various approaches that did not result in a performance improvement. We believe that reporting these negative results is just as important as highlighting our successes, as it provides a clearer picture of the technical challenges we faced and may help future researchers avoid similar pitfalls.

\subsection{Handcrafted Graph Features}

Before deciding on learned GNN embeddings, we explored the possibility of using explicit, handcrafted features derived from the heterogeneous graph.

\paragraph{Implemented Features}
We tested 11 additional features, including:
\begin{itemize}
    \item \textbf{Degree features}: Simple counts of connections for each identity type (device, IP, email, phone).
    \item \textbf{Historical fraud rates}: The fraud rate among other listings connected to the same identities. This was calculated using a point-in-time approach to ensure no data leakage occurred.
    \item \textbf{Recency}: The number of days since a particular identity was first seen on the platform.
    \item \textbf{Similarity features}: The fraud rate among listings that shared similar device characteristics but were not necessarily linked by the same fingerprint.
\end{itemize}

\paragraph{Performance Results}
The results of this experiment were disappointing. As shown in Table~\ref{tab:handcrafted_results}, adding these explicit features actually led to a significant drop in performance.

\begin{table}[h]
\centering
\caption{Performance degradation with handcrafted graph features}
\label{tab:handcrafted_results}
\begin{tabular}{lccc}
\hline
\textbf{Model} & \textbf{AUC-PR} & \textbf{Features} & \textbf{Change} \\
\hline
Vanilla XGBoost & 0.6639 & 65 & - \\
With Handcrafted Graph Features & 0.4833 & 76 (+11) & \textbf{-36.6\%} \\
\hline
\end{tabular}
\end{table}

\paragraph{Analysis of Failure}
A feature importance analysis revealed that these handcrafted features dominated the model's attention, overshadowing the more reliable base features. For instance, the IP-based fraud rate alone accounted for over 16% of the model's importance.

We believe there are three main reasons for this failure. First, there is a significant \textbf{distribution shift}: graph statistics computed during the training period did not generalize well to the test period because fraud patterns evolved too quickly. Second, these historical rates often introduced \textbf{spurious correlations} and noise that the model overfit to. Finally, the sheer weight of these new features caused a \textbf{"crowding out" effect}, where the model ignored the reliable signals from individual SEON attributes. We conclude that explicit graph features are inferior to both the vanilla XGBoost baseline and the learned GNN embeddings.

\subsection{Positive-Signal Edges Only}

We also tested whether filtering our graph to include only "high-risk" edges would improve the GNN's performance.

\paragraph{Hypothesis and Implementation}
The idea was that by restricting the graph to edges that correlate strongly with fraud (for example, device sharing specifically among known fraudsters), the GNN would learn more discriminative embeddings. We filtered the 31 million edges to retain only those where the connected nodes showed a fraud rate of at least 15%, which is roughly double the baseline.

\paragraph{Results and Analysis}
This approach provided no meaningful benefit, as shown in Table~\ref{tab:edge_filtering_results}.

\begin{table}[h]
\centering
\caption{Minimal benefit from positive-signal edge filtering}
\label{tab:edge_filtering_results}
\begin{tabular}{lcc}
\hline
\textbf{Configuration} & \textbf{AUC-PR} & \textbf{Change} \\
\hline
All edges (~31M edges) & 0.6990 & - \\
Positive-signal edges only (~8M edges) & 0.6967 & -0.3\% \\
\hline
\end{tabular}
\end{table}

It appears that supervised GNN training is already capable of weighting edges appropriately through backpropagation. Manual filtering likely removed potentially useful negative signals—such as legitimate users sharing IP addresses in public spaces—without providing any compensatory gain.

\subsection{Alternative Architectures: CARE-GNN}

We also evaluated the CARE-GNN (Camouflage-Resistant GNN) architecture \cite{dou2020enhancing}, which is specifically designed to identify fraudsters who attempt to mimic legitimate behavior.

\paragraph{Implementation Challenges}
Despite its theoretical advantages, CARE-GNN presented several practical challenges. It was significantly slower to train—taking 45 minutes per epoch compared to just 2 minutes for GraphSAGE—and required three times more memory due to its complex attention mechanisms. It was also highly sensitive to hyperparameters, particularly those related to label similarity.

\paragraph{Outcome and Conclusion}
Ultimately, CARE-GNN achieved an AUC-PR of 0.6945, which actually underperformed our GraphSAGE baseline. Given the 22-fold increase in training time and the lack of a performance gain, we concluded that GraphSAGE's simplicity and efficiency make it the better choice for this particular domain.

\subsection{Unsupervised GNN Training}

Finally, we looked at whether an unsupervised training approach could learn general graph structures that would be useful for fraud detection.

\paragraph{Hypothesis and Results}
We trained the GNN using unsupervised link prediction—essentially predicting whether an edge exists between two nodes—rather than supervised classification based on fraud labels. As Table~\ref{tab:unsupervised_results} shows, this method was significantly less effective.

\begin{table}[h]
\centering
\caption{Supervised GNN training vs. unsupervised}
\label{tab:unsupervised_results}
\begin{tabular}{lcc}
\hline
\textbf{GNN Training Method} & \textbf{AUC-PR} & \textbf{Notes} \\
\hline
Unsupervised link prediction & 0.6621 & Embeddings not fraud-specific \\
Supervised node classification & 0.6990 & \textbf{+5.6\% improvement} \\
\hline
\end{tabular}
\end{table}

The unsupervised embeddings captured the general topology of the network but failed to focus on the specific structural signatures of fraud. We learned that for this task, the direct signal from fraud labels is far more important than any general structural learning the unsupervised approach might provide.

\subsection{Summary of Alternative Approaches}

Our experiments with these alternative methods confirm that:
\begin{enumerate}
    \item Learned embeddings are far superior to handcrafted graph features.
    \item Supervised training is essential for GNN-based fraud detection.
    \item Simpler architectures like GraphSAGE often outperform more complex models when training efficiency is taken into account.
    \item Manual edge filtering adds complexity without a corresponding benefit in performance.
\end{enumerate}
